% \begin{itemize}
%     \item "Local permutation" of qubits (local, hence only within a single node) is allowed (i.e., a node can always relabel and reshuffle its own qubits among its memory slots)  
%     \item Graph states up to LC and local permutation (LP for short?). Together they form LCP (local complementation and permutation) equivalence classes
% \end{itemize}

We consider a set of network of nodes $\{A, B, \dots\}$, connected by quantum channels in an arbitrary topology.
Each node has a set of quantum memories; we denote by $Xi$ the qubit held by node $X$ in its $i^{th}$ memory.
Distributed multipartite states are graph states composed of those qubits.
%
The outcome of an entanglement distribution protocol is usually a graph state defined up to single-qubit Clifford corrections~\cite{}, 
which can be applied locally and are often considered inexpensive; furthermore, the memory label $i$ in which a node stores a qubit is arbitrary.
We thus abstract away these details and reason on \emph{equivalence classes} of graph states rather than on individual states.
First, in \cref{sec:lc} we abstract local Clifford operations and in \cref{sec:lp} we abstract local relabelings of qubits, then in \cref{sec:lcp} we define the equivalence class that we consider for our language in \cref{sec:semantics}.

\subsection{Local Cliffords (LC)}
\label{sec:lc}

\begin{wrapfigure}{r}{0.45\linewidth}
    \centering
    % ---------------------------------------------------------------------------
% LC equivalence: local complementation at A2 on the path A1-A2-A3.
%   N(A2) = {A1,A3} -> tau_{A2} toggles (A1,A3): path -> triangle.
% gold = edge added at this step; accent vertex = complemented vertex.
% ---------------------------------------------------------------------------
{\footnotesize
% panel: #1 = extra edges
\newcommand{\lcep}[1]{%
  \begin{tikzpicture}[baseline={(0,0.4)},x=1cm,y=1cm,
      gv/.append style={minimum size=6.5mm,font=\footnotesize}]
    \node[gv,frv] (a2) at (0,0.95) {$A2$};
    \node[gv] (a1) at (-0.75,0) {$A1$};
    \node[gv] (a3) at (0.75,0) {$A3$};
    \draw[ge](a1)--(a2);\draw[ge](a2)--(a3);
    #1
  \end{tikzpicture}}
\setlength{\tabcolsep}{4pt}%
\begin{tabular}{@{}c@{}c@{}c@{}}
\lcep{}
  & {\normalsize\tikz[baseline=-0.5ex]\draw[{To}-{To}] (0,0) --
      node[above,inner sep=2pt]{$\localComp{A2}$} (1.2,0);}
  & \lcep{\draw[newe](a1)--(a3);}
\end{tabular}%
}

    \caption{\acs{lc} equivalent graph states, up to local complementation on $A2$. 
    %\Acl{lc} $\localComp{A2}$ toggles the edge $A1 \edge A3$ (gold).
    }
    \label{fig:lc}
\end{wrapfigure}
We abstract the applications of local Clifford (LC) gates, 
as they are often considered \emph{free} operations in the literature~\cite{}; 
these are single-qubit gates, which typically have high fidelity~\cite{}. 
Ignoring local Cliffords is equivalent to ignoring local complementations (also conveniently shortened to LC) \alenote{Insert brief transition from stab states to graph states}
on the graph states, as local complementations are implemented by the Clifford unitary 
$\lcUnitary{} = \SX \bigotimesuNN S^\dagger$~\cite{vandennest2004graphical}. 
\acused{lc}
%
We hereafter refer to local Cliffords and local complementations interchangeably as simply \emph{\ac{lc}} operations.
%
We define the \emph{\ac{lc} equivalence} class as the set of states related by some sequence of \ac{lc} operations.
For instance, the two graph states illustrated in \cref{fig:lc} are \emph{\ac{lc} equivalent}.
% (up to \acl{lc} $\localComp{A2}$). 

We note that in some cases, a physical operation can produce non-LC-equivalent outputs from LC-equivalent inputs. 
For instance, as illustrated in \cref{fig:LC-input-nonLC-output}, two clearly LC-equivalent inputs (related by complementations $\localComp{A2} \localComp{B2} \localComp{C2}$) yield outputs in different LC equivalence classes after graph fusions $\fuse{A1,T1},\, \fuse{B1,T2},\, \fuse{C1,T3}$.
%
Albeit a non-trivial issue, this will not be problematic in our language as we discuss in \cref{sec:semantics}.

\begin{figure}[t]
    \centering
    \resizebox{0.9\textwidth}{!}
    {% ---------------------------------------------------------------------------
% LC-equivalent inputs, non-LC-equivalent outputs under graph fusion.
% Node T holds a triangle T1,T2,T3; nodes X in {A,B,C} each hold X1,X2,X3.
%   (a) X-parts are paths X1-X2-X3
%   (b) X-parts are triangles: tau_{X2} toggles (X1,X3).
%       So (b) = tau_{A2} tau_{B2} tau_{C2} (a).
% Fusions F(A1,T1), F(B1,T2), F(C1,T3) (green qubits, dashed pairing)
% toggle N(x) x N(t) and delete x,t:
%   (a) -> triangle A2,B2,C2 with pendants A3,B3,C3
%   (b) -> complete graph K6, LC-equivalent to GHZ_6
% The two outputs are not LC equivalent (checked by LC-orbit enumeration).
% Output vertices are drawn where they sit in the input.
% B- and C-parts point sideways, level with T2 and T3, to save height.
% gold = edge added by tau_{X2}.
% ---------------------------------------------------------------------------
{\footnotesize
% the X-part of node #1, rotated by #2 degrees around T and shifted
% vertically by #4; #3 = style of (X1,X3)
\newcommand{\lcfx}[4]{%
  \begin{scope}[shift={(0,#4)},rotate=#2]
    \node[gv,fuv] (#11) at (0,1.3) {$#11$};
    \node[gv] (#12) at (-0.42,1.95) {$#12$};
    \node[gv] (#13) at (0.42,1.95) {$#13$};
  \end{scope}
  \draw[ge](#11)--(#12);\draw[ge](#12)--(#13);\draw[#3](#11)--(#13);}
% input panel: #1 = style of the (X1,X3) edges (draw=none: absent)
\newcommand{\lcfin}[1]{%
  \begin{tikzpicture}[baseline={(0,0.5)},x=1cm,y=1cm,
      gv/.append style={minimum size=5.5mm,font=\scriptsize},
      fpair/.style={Ycol,line width=1.1pt,dash pattern=on 1.5pt off 1pt}]
    \node[gv,fuv] (T1) at (90:0.42)  {$T1$};
    \node[gv,fuv] (T2) at (210:0.42) {$T2$};
    \node[gv,fuv] (T3) at (330:0.42) {$T3$};
    \draw[ge](T1)--(T2)--(T3)--(T1);
    \lcfx{A}{0}{#1}{0}
    \lcfx{B}{90}{#1}{-0.21}
    \lcfx{C}{-90}{#1}{-0.21}
    \draw[fpair](A1)--(T1);\draw[fpair](B1)--(T2);\draw[fpair](C1)--(T3);
  \end{tikzpicture}}
% output panel: #1 = edges
\newcommand{\lcfout}[1]{%
  \begin{tikzpicture}[baseline={(0,0.5)},x=1cm,y=1cm,
      gv/.append style={minimum size=5.5mm,font=\scriptsize}]
    \foreach \X/\r/\s in {A/0/0,B/90/-0.21,C/-90/-0.21}{
      \begin{scope}[shift={(0,\s)},rotate=\r]
        \node[gv] (\X2) at (-0.42,1.95) {$\X2$};
        \node[gv] (\X3) at (0.42,1.95) {$\X3$};
      \end{scope}}
    \begin{pgfonlayer}{background}#1\end{pgfonlayer}
  \end{tikzpicture}}
% panel label, raised above the baseline
\newcommand{\lcflab}[1]{\raisebox{4mm}{\normalsize(#1)}}
\newcommand{\lcfarr}{{\small$\xrightarrow{\ \substack{\textstyle\fuse{A1,T1}\\[2pt]
    \textstyle\fuse{B1,T2}\\[2pt]\textstyle\fuse{C1,T3}}\ }$}}
\setlength{\tabcolsep}{1pt}%
\resizebox{0.85\linewidth}{!}{%
\begin{tabular}{@{}c@{}c@{}c@{}c@{}c@{}}
\lcflab{a} & \lcfin{draw=none} & \lcfarr
  & \lcfout{\draw[ge](A2)--(A3);\draw[ge](B2)--(B3);\draw[ge](C2)--(C3);
             \draw[ge](A2)--(B2)--(C2)--(A2);} & \lcflab{a'} \\\noalign{\vskip 12pt}
  & {\large$\Big\updownarrow\ \localComp{A2}\localComp{B2}\localComp{C2}$}
  & & {\Large$\not\sim_{\mathrm{LC}}$} \\\noalign{\vskip 12pt}
\lcflab{b} & \lcfin{newe} & \lcfarr
  & \lcfout{\foreach \u/\v in {A2/A3,A2/B2,A2/B3,A2/C2,A2/C3,A3/B2,A3/B3,
               A3/C2,A3/C3,B2/B3,B2/C2,B2/C3,B3/C2,B3/C3,C2/C3}
               {\draw[ge](\u)--(\v);}} & \lcflab{b'}
\end{tabular}}%
}
}
    \caption{\Ac{lc} equivalent inputs yielding non-\ac{lc}-equivalent outputs. 
    The inputs (a) and (b) differ by $\localComp{A2}\localComp{B2}\localComp{C2}$, 
    which adds the gold edges in (b). 
    Fusing the green qubits paired by green dashed lines yields a triangle with three leaves 
    (a') and the complete 6-qubit graph (b'), which are not LC~equivalent.}
    \label{fig:LC-input-nonLC-output}
\end{figure}

\subsection{Local Permutation (LP)}
\label{sec:lp}

We abstract the physical locations of qubits at a network node by allowing local relabeling of qubits.
For instance, we want the two graph states of \cref{fig:lp} 
%$A1 \edge A2 \edge A3 \edge A4$ and $A1 \edge A3 \edge A2 \edge A4$,  i.e., states $\ket{0000} + \ket{0011} + \ket{1100} - \ket{1111}$ and $\ket{0000} + \ket{0101} + \ket{1010} - \ket{1111}$, 
to be considered equivalent, as the labels naming the qubits are arbitrary; both graph states are linear graphs of four qubits. We define \emph{\ac{lp}} such a relabeling of qubits. 
%
We only allow \emph{local} permutations, i.e., within the same network node, as permutations of distributed qubits would require (lossy) transmission of said qubits or non-local operations. % (e.g., a SWAP gate).
%
An \ac{lp} defines\ale{is?} a permutation matrix $P$ that maps the adjacency matrix $\Gamma$ of a graph to $P \Gamma P^\top$; 
$P$ only permutes qubits of the same node, thus it is block diagonal. 
For instance, a relabeling exchanging $\permnoletter{A1}{A2}$ in \cref{fig:lp} corresponds to the matrix shown below the arrow.

\begin{wrapfigure}{r}{0.52\linewidth}
    \centering
    % All qubits are at node A; labels are arbitrary within the node.
\begin{tikzpicture}[baseline=(current bounding box.center),
    gv/.append style={minimum size=5.5mm,font=\small}]
    \foreach \i/\lab in {0/A1,1/A2,2/A3,3/A4}{
        \node[gv] (l\i) at (.85*\i,0) {$\lab$};
    }
    \draw[ge] (l0)--(l1)--(l2)--(l3);
    \draw[bigarr] (3.1,0) -- node[above,font=\small] {$\permnoletter{A2}{A3}$} (4.5,0);
    \foreach \i/\lab in {0/A1,1/A3,2/A2,3/A4}{
        \node[gv] (r\i) at (5.05+.85*\i,0) {$\lab$};
    }
    \draw[ge] (r0)--(r1)--(r2)--(r3);
\end{tikzpicture}

    \caption{\Ac{lp}-equivalent graph states, up to local permutation $\permnoletter{A2}{A3}$. 
    }
    \label{fig:lp}
\end{wrapfigure}
%
We define the \emph{\ac{lp} equivalence} class as the set of states related by some \ac{lp}.
We note that \ac{lp} equivalence and \ac{lc} equivalence are not correlated, as far as we have observed.\alenote{Validate/Check this?}
For instance, the two graph states illustrated in \cref{fig:lp} are \emph{\ac{lp} equivalent}, but not \ac{lc} equivalent, while the states in \cref{fig:lc} are \ac{lc} but not \ac{lp} equivalent 

\subsection{Local Complementation and Permutation (LCP)}
\label{sec:lcp}

An interesting observation arises when considering LC and LP equivalence together. 
We define the \emph{\ac{lcp} equivalence} class as the set of states related by some sequence of \ac{lc} operations and \acp{lp}.
We denote by $\lcpClass{G}$ the \ac{lcp} class of a graph state $\ket{G}$, and by $\LcpSet$ the set of all \ac{lcp} classes.
For every class in $\LcpSet$, we fix a \emph{canonical representative}, \ale{Specify how the canonical representative is defined? Or don't mention this?} \alenote{We may not need to specify a canonical rep of LCP classes. We could define them as "the LCP eq. class of G is the set of graph states reachable by applying any LP or LC to G". Then in the implementation we will use the canonical rep for practicality of establishing LCP equivalence.}

We note that \acp{lc} and \acp{lp} commute simply as
$$
\localComp{A1}\perm{A1}{A2} = \perm{A1}{A2}\localComp{A2}.
$$
Thus, \acs{lcp} equivalence can also be defined as the set of states related by one \acs{lp} followed by any sequence of \ac{lc} operations.
% More simply, we can define \acs{lcp} equivalence as one \acs{lp} relating two \ac{lc} equivalence classes.
From this definition, we observe that \acp{lp} relate different whole \ac{lc} equivalence classes.  
% Thus, \ac{lcp} equivalence classes strictly contain entire \ac{lc} equivalence classes.
%
To show that multiple \ac{lc} equivalence classes can be contained in a single \ac{lcp} class, \cref{fig:lcp} illustrates two graph states that are \ac{lcp} equivalent but neither \ac{lc} nor \ac{lp} equivalent.

\begin{figure}
    \centering
    % ---------------------------------------------------------------------------
% LCP equivalence: path B1-A1-A2-B2  --P(A1<->A2)-->  path B1-A2-A1-B2
%   --tau_{A1}-->  N(A1)={A2,B2}, toggles (A2,B2).
% Left and right graphs are LCP equivalent, but neither LC equivalent
% (checked by LC-orbit enumeration) nor LP equivalent (not isomorphic).
% Same graphs as lp_equiv.tex up to renaming; dashed boxes = network nodes
% A (top) and B (bottom).  Vertices drawn in the path order of the middle panel.
% gold = edge added at this step; accent vertex = complemented vertex.
% ---------------------------------------------------------------------------
{\footnotesize
% panel: #1..#4 = labels in path order ; #5 = style of 3rd vertex ;
%        #6 = extra edges
\newcommand{\lcpp}[6]{%
  \begin{tikzpicture}[baseline={(0,0.4)},x=1cm,y=1cm,
      gv/.append style={minimum size=6.5mm,font=\footnotesize},
      netn/.style={draw=Vcol!70,dashed,rounded corners=4pt,fill=##1,
                   fill opacity=0.12,inner sep=2.5pt}]
    \node[gv] (p1) at (0,0) {$#1$};
    \node[gv] (p2) at (0,1.05) {$#2$};
    \node[gv,#5] (p3) at (1.05,1.05) {$#3$};
    \node[gv] (p4) at (1.05,0) {$#4$};
    \begin{pgfonlayer}{background}
    \node[netn=Zcol,fit=(p2)(p3)] (nA) {};
    \node[netn=Ycol,fit=(p1)(p4)] (nB) {};
    \draw[ge](p1)--(p2)--(p3)--(p4);
    #6
    \end{pgfonlayer}
    \node[font=\scriptsize\bfseries,text=Vcol,anchor=east] at (nA.west) {A};
    \node[font=\scriptsize\bfseries,text=Vcol,anchor=east] at (nB.west) {B};
  \end{tikzpicture}}
\setlength{\tabcolsep}{2pt}%
\resizebox{0.75\linewidth}{!}{%
\begin{tabular}{@{}c@{\hspace{6pt}}c@{\hspace{6pt}}c@{\hspace{6pt}}c@{\hspace{6pt}}c@{}}
\lcpp{B1}{A1}{A2}{B2}{}{}
  & {\normalsize\tikz[baseline=-0.5ex]\draw[{To}-{To}] (0,0) --
      node[above,inner sep=2pt]{\small$\permnoletter{A1}{A2}$} (1.9,0);}
  & \lcpp{B1}{A2}{A1}{B2}{frv}{}
  & {\normalsize\tikz[baseline=-0.5ex]\draw[{To}-{To}] (0,0) --
      node[above,inner sep=2pt]{$\localComp{A1}$} (1.0,0);}
  & \lcpp{B1}{A2}{A1}{B2}{frv}{\draw[newe](p2)--(p4);}
\end{tabular}}%
}

    \caption{\Ac{lcp} equivalent graph states (left and right), related by $\perm{A1}{A2}$ followed by $\localComp{A1}$, which are neither \ac{lc} nor \ac{lp} equivalent.}
    \label{fig:lcp}
\end{figure}